\documentclass[11pt]{amsart}
\usepackage[letterpaper,margin=1in]{geometry}
\usepackage{amsmath,amssymb,amsthm,mathtools,amscd}
\usepackage{microtype}
\usepackage[hidelinks]{hyperref}
\usepackage{enumitem}
\usepackage{booktabs,longtable,array}
\usepackage{fontspec}
\usepackage{unicode-math}
\setmainfont{Times New Roman}
\setsansfont{Arial}
\setmonofont{Menlo}[Scale=MatchLowercase]
\setmathfont{texgyretermes-math.otf}
\newfontfamily\AgdaFont{Arial Unicode MS}[Scale=MatchLowercase]
\def\UrlFont{\AgdaFont}
\setlength{\parskip}{0.25em}
\setlength{\parindent}{1.25em}
\setlength{\emergencystretch}{2em}
\numberwithin{equation}{section}
\newtheorem{definition}{Definition}[section]
\newtheorem{example}[definition]{Example}
\newtheorem{lemma}[definition]{Lemma}
\newtheorem{proposition}[definition]{Proposition}
\newtheorem{theorem}[definition]{Theorem}
\newtheorem{corollary}[definition]{Corollary}
\theoremstyle{remark}
\newtheorem{remark}[definition]{Remark}
\newtheorem{convention}[definition]{Convention}
\newtheorem{warning}[definition]{Warning}
\newcommand{\BG}{\mathbf{B}G}
\newcommand{\GTors}{G\mathrm{Tors}_{L}}
\newcommand{\GTorsLevel}{G\mathrm{Tors}_{L,\ell}}
\newcommand{\Loop}{\mathop{\Omega}}
\newcommand{\Trunc}[1]{\lVert #1 \rVert}
\newcommand{\Type}{\mathsf{Type}}
\newcommand{\agda}[1]{{\AgdaFont #1}}
\newcounter{formalization}[section]
\renewcommand{\theformalization}{\thesection.\arabic{formalization}}
\newcommand{\formal}[3]{\par\smallskip\noindent\begingroup\footnotesize\sffamily\setlength{\parindent}{0pt}\leftskip=1em\rightskip=0pt plus 1fil\refstepcounter{formalization}\textbf{Formalization~\theformalization.} \textbf{File:} #1. \textbf{Declaration:} \agda{#2}. #3\par\endgroup\smallskip}
\newcommand{\upstream}[2]{\par\smallskip\noindent\begingroup\footnotesize\sffamily\setlength{\parindent}{0pt}\leftskip=1em\rightskip=0pt plus 1fil\refstepcounter{formalization}\textbf{Upstream interface~\theformalization.} \textbf{File:} #1. \textbf{Declaration:} \agda{#2}.\par\endgroup\smallskip}
% Retain amsart's article layout while matching the title block used by the
% GroupPresentations article in this repository.
\makeatletter
\def\@settitle{%
  \begin{center}
    \LARGE\bfseries \@title
  \end{center}}
\def\@setauthors{%
  \begin{center}
    \large
    \def\\{\protect\linebreak}%
    \authors
  \end{center}}
\makeatother
\title[Torsors, Univalence, and Delooping]{Torsors, Univalence, and Delooping\\[0.35em]\large\textnormal{A Mathematical Companion to a Cubical Agda Formalization}}
\author{Hongwei (Alex) Wang\\[3pt]
\small New York University Abu Dhabi\\
\small \texttt{hw4375@nyu.edu}}
\dedicatory{\normalfont\normalsize 14 November 2026}

\begin{document}
\begin{abstract}
We give a mathematical account of a Cubical Agda development of generic left torsors and their moduli. The article develops the regular torsor, univalent paths between torsors, the connected $1$-type of torsor moduli, and its comparison with the generic Eilenberg--Mac Lane delooping. It also proves a universe-lifted path-torsor characterization under explicit pointed loop-group data. Compact annotations identify the checked local declarations and imported Cubical interfaces; they are provenance for, rather than replacements for, the mathematical arguments. The result is an equivalence of underlying types at the stated universe levels, not a separately packaged group-isomorphism theorem or a construction of computational transport.
\end{abstract}
\maketitle
\pagestyle{plain}

\section{Why Torsors?}
\label{ch:why}

Symmetry is often presented by a group acting on a set. The group records transformations and its unit is a distinguished transformation. In many geometric and physical situations, however, the object on which the transformations act has no preferred origin. A coordinate frame in an affine space, a choice of phase for a nonzero complex number, and a choice of trivialization of a principal bundle are all examples: after one choice has been made, every other choice is obtained by a unique symmetry, but no choice is privileged.

The correct object for this situation is a torsor. A torsor remembers the action of a group but deliberately forgets the identity element as a point of the acted-on space. The loss of a chosen point is not a defect. It is precisely what makes torsors appropriate for moduli problems and for the construction of classifying spaces.

\subsection{The regular action}
Let $G$ be a group, written multiplicatively, with unit $1$. The most familiar left action of $G$ is its action on itself by left multiplication:
\begin{equation}
g\cdot x \coloneqq gx.
\end{equation}
The action laws are $1\cdot x=x$ and $(gh)\cdot x=g\cdot(h\cdot x)$. The carrier $G$ does have a preferred point, namely $1$, but that point is extra data: the action itself does not mention it.

\begin{example}[Changing origins]
Let $V$ be a vector space. The translations of $V$ form an additive group. An affine space modelled on $V$ has points but no canonical zero. Given two affine points $x$ and $y$, there is a unique translation carrying $x$ to $y$. Thus an affine space is a torsor for the translation group. Choosing an origin identifies it with $V$, but different origins give different identifications.
\end{example}

\begin{example}[Phases]
The unit circle group $U(1)$ acts on a one-dimensional complex Hilbert space by scalar multiplication. On the set of unit vectors in a fixed complex line, this action is free and transitive. A unit vector is a phase choice; the torsor records phase differences without selecting a reference phase.
\end{example}

\subsection{Torsors as symmetry without an origin}
The phrase ``free and transitive'' is traditionally used for torsors. Freeness says that no nonidentity group element fixes a point. Transitivity says that any point can be reached from any other. Together they say that, once an origin $x$ is temporarily selected, the orbit map
\begin{equation}
\operatorname{orb}_x:G\longrightarrow X, \qquad g\longmapsto g\cdot x
\end{equation}
is a bijection. In homotopy type theory, ``bijection'' is replaced by equivalence, and the most useful formulation says that every fiber of the orbit map is contractible. This retains the unique-transporter idea while behaving correctly in a homotopical setting.

\begin{remark}[Why ``mere'' inhabitation?]
A torsor should not contain a selected point. Therefore one asks only that its carrier be inhabited after propositional truncation, written $\Trunc{X}_{-1}$, rather than providing a point of $X$. This supplies enough existence to prove invariant propositions, but it prevents a construction from extracting a canonical origin. The distinction is essential in the connectedness proof for the torsor moduli type.
\end{remark}

\subsection{From sets to types}
The present development works with a set-level group. Its carrier is a set in the homotopy-theoretic sense, so equality between group elements is a proposition. A torsor carrier is also required to be a set. Consequently, the interesting higher structure appears not inside an individual torsor but in the type of all torsors: paths between torsors encode equivalences of structures.

The central theme of this article is that this moduli type behaves as a delooping of $G$. Informally, a delooping $\BG$ is a connected pointed $1$-type whose loop space at the basepoint has underlying type equivalent to $G$. The word ``underlying'' matters. An equivalence of loop types is not automatically a group isomorphism with respect to a chosen multiplication unless compatibility is proved and packaged.

\subsection{Formal scope}
The local Cubical Agda development begins with generic groups and generic torsors. It does not yet specialize to $\mathbb{C}^{\times}$, construct geometric principal bundles, or define computational transport. Those are later applications. The present result is a reusable homotopy-theoretic foundation.

\formal{\texttt{Regular.agda}}{TorsorDelooping.Regular.regularLeftAction; regularLeftTorsor}{The regular left action and its torsor structure are formalized for every set-level Cubical group.}
\upstream{\texttt{Properties.agda}}{·GroupAutomorphismR}

\section{Left Torsors and Equivariant Maps}
\label{ch:left-torsors}

\subsection{Actions and transporters}
Fix a set-level group $G$ and a set $X$. A left action is a function
\begin{equation}
\alpha:G\times X\longrightarrow X, \qquad (g,x)\longmapsto g\cdot x,
\end{equation}
satisfying the two usual laws. We use left actions throughout. Thus the order convention is
\begin{equation}
(gh)\cdot x=g\cdot(h\cdot x).
\label{eq:left-action-law}
\end{equation}
Keeping this convention visible prevents a later ambiguity: automorphisms of the regular left torsor are naturally given by \emph{right} multiplication.

\begin{definition}[Transporter and left torsor]
For $x,y:X$, define the transporter type
\begin{equation}
\operatorname{Trans}_G(x,y)\coloneqq\sum_{g:G}(g\cdot x=y).
\end{equation}
A left $G$-torsor is a set $X$ with a left action, a proof of $\Trunc{X}_{-1}$, and a proof that $\operatorname{Trans}_G(x,y)$ is contractible for all $x,y:X$.
\end{definition}

Contractibility says more than nonemptiness. It gives a centre $(g,p)$ and says that every other group element/path pair is equal to it. When all types involved are sets, this is exactly the assertion that there is one and only one element transporting $x$ to $y$.

\formal{\texttt{LeftAction.agda}}{TorsorDelooping.LeftAction.LeftAction}{The fields are \texttt{isSetX}, \texttt{act}, \texttt{act-unit}, and \texttt{act-comp}; Equation \eqref{eq:left-action-law} is the last field.}
\formal{\texttt{LeftTorsor.agda}}{TorsorDelooping.LeftTorsor.IsLeftTorsor; LeftTorsor}{Mere inhabitation is the upstream \agda{∥ X ∥₁}, Cubical's notation for propositional (hence $(-1)$-) truncation; the contractible transporter is the field \texttt{regular}.}

\begin{proposition}[Orbit criterion]
Let $X$ carry a left $G$-action. The following are equivalent.
\begin{enumerate}[label=\textup{(\roman*)}]
\item For every $x,y:X$, $\operatorname{Trans}_G(x,y)$ is contractible.
\item For every $x:X$, the orbit map $\operatorname{orb}_x:G\to X$, $g\mapsto g\cdot x$, is an equivalence.
\end{enumerate}
\end{proposition}
\begin{proof}
The fiber of $\operatorname{orb}_x$ over $y$ is, by definition,
$\sum_{g:G}(g\cdot x=y)$. Thus condition (i) says exactly that every fiber is contractible, which is the standard characterization of an equivalence. Conversely an equivalence has contractible fibers, giving (i).
\end{proof}

\begin{remark}
In the local development the orbit criterion is used in the direction from contractible transporters to equivalences. It is named \texttt{orbitEquiv}. The proof relies only on the definition of \texttt{isEquiv} by contractible fibers.
\end{remark}

\formal{\texttt{Morphisms.agda}}{TorsorDelooping.Morphisms.orbitEquiv}{The checked result sends a torsor, a point $x$, and a target $y$ to the equivalence induced by the corresponding contractible transporter.}

\subsection{The regular torsor}
\begin{proposition}[The regular left torsor]
The left action of $G$ on itself is a left $G$-torsor.
\end{proposition}
\begin{proof}
Mere inhabitation is witnessed by $1\in G$. For $x,y:G$, a transporter is an element $g$ together with $gx=y$. The unique solution is $g=yx^{-1}$. Indeed, $(yx^{-1})x=y$ by associativity and inverse cancellation. If $gx=y$, multiplying on the right by $x^{-1}$ gives $g=yx^{-1}$; all equality data are propositions because $G$ is a set. Hence the transporter is contractible.
\end{proof}

Notice a subtle but useful asymmetry. The action is by left multiplication, whereas the proof may equivalently identify the transporter by the right translation $g\mapsto gx$. Cubical's upstream group library supplies this right-translation automorphism directly.

\formal{\texttt{Regular.agda}}{TorsorDelooping.Regular.regularLeftTorsor}{The centre of mere inhabitation is the unit. Regularity is discharged through the upstream right-multiplication automorphism.}

\subsection{Equivariant maps}
\begin{definition}
Let $X$ and $Y$ be left $G$-actions. A function $f:X\to Y$ is equivariant when
\begin{equation}
f(g\cdot x)=g\cdot f(x)
\end{equation}
for all $g$ and $x$. Equivariant maps have identities and compose.
\end{definition}

\begin{proposition}[An equivariant map of torsors is an equivalence]
Let $f:X\to Y$ be an equivariant map between left $G$-torsors. Then $f$ is an equivalence.
\end{proposition}
\begin{proof}
The assertion that $f$ is an equivalence is a proposition. We may therefore eliminate the merely inhabited source torsor and temporarily choose $x:X$. Equivariance says precisely that the two orbit maps satisfy
\begin{equation}
  f \circ \operatorname{orb}_x =
  \operatorname{orb}_{f(x)}.
\end{equation}
The two orbit maps are equivalences by the orbit criterion. Hence $f$ is a composite of equivalences, and therefore an equivalence. Since the conclusion was propositional, the temporary choice of $x$ is eliminated.
\end{proof}

The displayed equality is only mnemonic; no commutative-square package is needed in the checked proof. The proof is nevertheless significant: it explains why morphisms in the groupoid of torsors are automatically invertible.

\formal{\texttt{Morphisms.agda}}{TorsorDelooping.Morphisms.equivariant-id; equivariant-comp; equivariantMapIsEquiv}{The final proof uses propositional-truncation elimination into the proposition \texttt{isEquiv f}.}

\subsection{Pointings and trivializations}
A point $x:X$ of a torsor is a trivialization: the orbit equivalence $G\simeq X$ identifies $g$ with $g\cdot x$. A different point changes this identification by a translation. Thus pointed torsors are equivalent to copies of $G$, while unpointed torsors are the correct origin-free objects. This distinction will reappear when the regular torsor is used as a basepoint of a moduli type.

\section{The Moduli Type of Torsors}
\label{ch:moduli}

\subsection{Why form a moduli type?}
A single torsor records a symmetry situation without a preferred origin. The type of all torsors records how such situations vary. In ordinary mathematics one often calls the resulting groupoid the groupoid of $G$-torsors: objects are torsors and arrows are equivariant equivalences. Homotopy type theory presents the same groupoid internally, with equalities of structured objects playing the role of isomorphisms.

This change of viewpoint is decisive. Instead of declaring isomorphic torsors equivalent by an external convention, univalence identifies equivalence of carriers with equality of types, and the Structure Identity Principle extends that identification to the torsor structures.

\subsection{Structure on a carrier}
Fix a universe level $\ell$. For a set $X: \Type_\ell$, the data of an action comprise an operation $G\to X\to X$ and its unit and associativity laws. The torsor axioms add mere inhabitation and contractible transporters. Write
\begin{equation}
\GTorsLevel
\end{equation}
for the type of a carrier equipped with this full structure. In the checked development this type lies in $\Type_{\ell+1}$ when $G: \Type_\ell$: quantifying over carriers in $\Type_\ell$ raises the universe level.

\begin{remark}[Why the axioms are propositions]
The law fields of an action are propositions because the carrier is a set. Mere inhabitation is a proposition by construction, and contractibility is a proposition. Therefore the structure has no further choices beyond its operation. This is exactly the hypothesis that allows a univalent structure theorem to apply.
\end{remark}

\formal{\texttt{Moduli.agda}}{TorsorDelooping.Moduli.ActionAxioms; TorsorAxioms; isPropActionAxioms; isPropTorsorAxioms}{The operation is inherited from Cubical's left-action structure; all additional law fields are proved propositional.}

\subsection{Equivalences of torsors}
Given torsors $T$ and $U$, an equivariant equivalence is an equivalence of their carriers which commutes with the actions. Write
\begin{equation}
\operatorname{TorsEquiv}(T,U)
\end{equation}
for this type. It is important that the equivariance proof is part of the structure. It is not enough merely to know that the carriers are equivalent: arbitrary equivalences need not respect the $G$-action.

\begin{definition}[Structure identity for torsors]
The Structure Identity Principle for torsors is an equivalence
\begin{equation}
\operatorname{TorsEquiv}(T,U) \simeq (T=U).
\label{eq:sip}
\end{equation}
Thus a path in the moduli type is represented by an equivariant equivalence, and every equivariant equivalence produces such a path.
\end{definition}

\paragraph{Mathematical explanation.}
Univalence identifies an equivalence $e:X\simeq Y$ with a path $X=Y$ in the universe. Transporting the action structure along that path asks precisely that $e$ conjugate the action on $X$ to the action on $Y$. The remaining torsor fields are propositions, so their transports carry no independent choices. Consequently, an equivalence of carriers plus equivariance is exactly a path of torsor structures. The formal proof applies the library SIP theorem to the nested action and torsor structures.

\formal{\texttt{Moduli.agda}}{TorsorDelooping.Moduli.torsorPathEquiv}{This is the exact checked form of Equation \eqref{eq:sip}.}
\upstream{\texttt{SIP.agda}}{SIP}

\subsection{What equality now means}
Equation \eqref{eq:sip} is not a slogan that ``isomorphic objects are equal'' in an unqualified sense. It is an equivalence between an explicitly specified type of structured equivalences and an identity type of a particular univalent structure. It depends on the carrier being in a universe, on univalence for that universe, and on the proof-irrelevance of the axiom fields.

\begin{example}[A torsor and its translate]
Let $X$ be a $G$-torsor and let $a:X$. The orbit equivalence $g\mapsto g\cdot a$ identifies the regular torsor with $X$. By the SIP it gives a path from the regular point of the moduli type to the point represented by $X$. Replacing $a$ by another point generally gives a different path, but the existence of such a path is all that connectedness requires.
\end{example}

\subsection{A warning about paths and composition}
An equality in a moduli type has a concatenation operation. An equivariant equivalence has composition. Equation \eqref{eq:sip} alone does not permit one to silently replace concatenation by composition: a coherence theorem is required, and its orientation depends on conventions. The regular-torsor coherence needed later is proved separately using the computation of SIP on carrier paths and the univalence law for composition.

\begin{warning}
The paper will use the SIP to pass between paths and equivariant equivalences, but it will not assert a general unstated theorem that SIP preserves composition. Only the restricted regular-translation coherence actually checked in the development is used for the $\mathsf{EM}_1$ comparison.
\end{warning}

\section{Torsor Moduli as a Delooping}
\label{ch:delooping}

\subsection{The regular basepoint}
The regular left torsor defines a canonical point
\begin{equation}
b_G:\mathbf{1}\longrightarrow \GTorsLevel.
\end{equation}
We call its value the regular basepoint. Although a general torsor has no selected element, the regular torsor has a canonical element $1\in G$, so it is the natural reference object in the moduli type.

\formal{\texttt{Delooping.agda}}{TorsorDelooping.Delooping.torsorBase}{The basepoint is obtained by embedding \texttt{regularLeftTorsor G} into the SIP moduli type.}

\subsection{Connectedness}
\begin{proposition}
The moduli type $\GTorsLevel$ is connected.
\end{proposition}
\begin{proof}
Let $T$ be a torsor. Its carrier is merely inhabited. To prove that the regular basepoint and $T$ are merely path equal, we may eliminate this truncation into the proposition $\Trunc{b_G=T}_{-1}$. Thus choose $x:T$ temporarily. The orbit map
\begin{equation}
G\longrightarrow T, \qquad g\longmapsto g\cdot x,
\end{equation}
is equivariant for the regular action and the action of $T$. By Proposition 2.2 it is an equivalence; by the SIP it yields a path $b_G=T$. Reintroducing propositional truncation gives a mere path. Since every point is merely connected to $b_G$, the moduli type is connected.
\end{proof}

The proof illustrates a characteristic use of propositional truncation. We choose a point only while constructing a proposition, so the result does not select an origin for any torsor.

\formal{\texttt{Delooping.agda}}{TorsorDelooping.Delooping.merelyBasePath; isConnectedLeftTorsorModuli}{The implementation constructs the equivariant orbit map, converts it to a SIP path, and eliminates mere inhabitation into a mere path.}

\subsection{The moduli is a one-type}
\begin{proposition}
The moduli type $\GTorsLevel$ is a $1$-type.
\end{proposition}
\begin{proof}
Given torsors $T,U$, Equation \eqref{eq:sip} identifies their path type with the type of equivariant equivalences $T\simeq U$. The carrier types are sets. The type of equivalences between sets is a set, and the equivariance condition is a proposition. Therefore the type of equivariant equivalences is a set. Hence every path type of the moduli is a set, which is exactly the assertion that the moduli is a $1$-type.
\end{proof}

\formal{\texttt{Delooping.agda}}{TorsorDelooping.Delooping.isGroupoidLeftTorsorModuli}{Cubical's name \texttt{isGroupoid} expresses that all path types are sets, i.e. that the type is $1$-truncated.}

\subsection{Automorphisms of the regular torsor}
Loops at $b_G$ are, by the SIP, equivariant automorphisms of the regular torsor. Such an automorphism has a concrete form.

\begin{proposition}[Regular-torsor automorphisms]
Evaluation at the identity induces an equivalence
\begin{equation}
\operatorname{Aut}_G(G)\simeq G,
\end{equation}
where $\operatorname{Aut}_G(G)$ denotes the equivariant self-equivalences of the regular left torsor. An element $h:G$ corresponds to right multiplication
\begin{equation}
R_h(x)\coloneqq xh.
\label{eq:right-translation}
\end{equation}
\end{proposition}
\begin{proof}
Right multiplication is equivariant because
\begin{equation}
R_h(gx)=(gx)h=g(xh)=gR_h(x).
\end{equation}
It is invertible with inverse $R_{h^{-1}}$. Conversely, if $f:G\to G$ is equivariant, then
\begin{equation}
f(x)=f(x\cdot 1)=x\cdot f(1)=xf(1).
\end{equation}
Thus $f=R_{f(1)}$. Evaluation at $1$ and $h\mapsto R_h$ are mutually inverse.
\end{proof}

\formal{\texttt{Delooping.agda}}{TorsorDelooping.Delooping.regularAutIsoGroup}{The inverse map is \texttt{rightRegularAut}; the forward map evaluates at the group unit.}

\subsection{Loop types and orientation}
Combining the inverse SIP equivalence with the preceding proposition gives
\begin{equation}
\Loop(\GTorsLevel,b_G)\simeq G
\label{eq:looptype}
\end{equation}
as an equivalence of underlying types. Its direction in the source code is from loops to the carrier of $G$.

\formal{\texttt{Delooping.agda}}{torsorLoopEquivGroup}{The checked declaration gives the underlying equivalence from the basepoint loop type to the carrier of $G$.}

\begin{warning}[Type equivalence is not yet a group isomorphism]
Loop concatenation equips the left side of \eqref{eq:looptype} with a group structure. The displayed checked declaration alone does not package preservation of multiplication, units, and inverses. The next section proves the particular multiplication coherence needed for the $\mathsf{EM}_1$ recursor; it should not be read as a separately packaged general loop-group isomorphism theorem.
\end{warning}

\subsection{The SIP composition coherence actually used}
Let $p(g)$ denote the loop obtained from the right-translation automorphism $R_g$ by the SIP. With the convention that composition first applies $R_g$ and then $R_h$, Equation \eqref{eq:right-translation} gives
\begin{equation}
R_h(R_g(x))=(xg)h=x(gh)=R_{gh}(x).
\end{equation}
The local composition convention therefore yields an equality of equivariant equivalences corresponding to $R_{gh}$. Passing to paths requires an additional calculation.

\begin{proposition}[Regular SIP coherence]
For $g,h:G$,
\begin{equation}
p(gh)=p(g)\mathbin{\cdot}p(h),
\label{eq:sip-coherence}
\end{equation}
where the product on the right is path concatenation.
\end{proposition}
\begin{proof}
The carrier path underlying the SIP image of an equivariant equivalence is its univalence path. The univalence composition theorem identifies the carrier path of the composite $R_g$ followed by $R_h$ with the concatenation of the carrier paths of $R_g$ and $R_h$. The two structured paths have equal carrier paths. Since the remaining structure fields are propositional, equality of their carrier paths lifts to equality of the full torsor paths. Finally the equality of right translations identifies the composite with $R_{gh}$. This proves \eqref{eq:sip-coherence}.
\end{proof}

\formal{\texttt{LoopGroup.agda}; \texttt{SIPPathEmbedding.agda}; \texttt{SIPCoherence.agda}}{rightRegularAut-comp; caracTorsor≡; regularTorsorLoop-comp}{The proof combines the local right-translation composition law, the structured-path embedding, and the upstream univalence-composition theorem.}
\upstream{\texttt{Univalence.agda}}{uaCompEquiv}

\section{Comparison with \texorpdfstring{$\mathsf{EM}_1$}{EM1}}
\label{ch:em1}

\subsection{The Eilenberg--Mac Lane delooping}
For a set-level group $G$, Cubical Agda supplies a higher inductive type $\mathsf{EM}_1(G)$. It has a basepoint and one generating loop for every $g:G$, together with constructors imposing the group multiplication law and a truncation constructor making the result a $1$-type. This is a synthetic model of the usual classifying space of a discrete group.

The relevant recursor has a particularly transparent form. To define a map from $\mathsf{EM}_1(G)$ into a $1$-type $B$, it suffices to give:
\begin{enumerate}[label=\textup{(\roman*)}]
\item a point $b:B$;
\item a loop $\lambda_g:b=b$ for each $g:G$; and
\item paths $\lambda_{gh}=\lambda_g\mathbin{\cdot}\lambda_h$.
\end{enumerate}
The target's $1$-truncation condition is what makes this data sufficient.

\upstream{\texttt{Properties.agda}}{rec'}

\subsection{The comparison map}
Set $B=\GTorsLevel$ and take $b=b_G$. For the loop indexed by $g$, use the regular-torsor loop $p(g)$. The coherence required by the recursor is exactly Proposition 4.5. Thus one obtains
\begin{equation}
f:\mathsf{EM}_1(G)\longrightarrow \GTorsLevel.
\label{eq:em-map}
\end{equation}
By construction $f$ sends the Eilenberg--Mac Lane basepoint to the regular torsor and sends the generator loop for $g$ to $p(g)$.

\formal{\texttt{EM1Comparison.agda}}{TorsorDelooping.EM1Comparison.em1ToTorsor; em1ToTorsor-base; em1ToTorsor-loop}{The map is defined by upstream \texttt{rec'} with target \texttt{isGroupoidLeftTorsorModuli G}, basepoint \texttt{torsorBase G}, and coherence \texttt{regularTorsorLoop-comp G}.}

\subsection{The basepoint loop map}
The upstream library proves that the loop type of $\mathsf{EM}_1(G)$ at its basepoint is equivalent to the carrier of $G$. The local development proves that $p:G\to \Loop(\GTorsLevel,b_G)$ is an equivalence. Consequently, the map on basepoint loops induced by $f$ is an equivalence.

\begin{proposition}
The pointed map in \eqref{eq:em-map} induces an equivalence on loops at the displayed basepoints.
\end{proposition}
\begin{proof}
Identify a loop in $\mathsf{EM}_1(G)$ with an element $g:G$ using the upstream equivalence $\Omega\mathsf{EM}_1\mathsf{Iso}$. The computation rule for $f$ identifies the image of the loop with $p(g)$. The map $p$ is an equivalence by the regular-torsor loop calculation. Their composite is therefore an equivalence.
\end{proof}

\formal{\texttt{EM1Comparison.agda}}{TorsorDelooping.EM1Comparison.isEquiv-em1ToTorsor-baseLoop}{The proof combines the upstream loop calculation with the local regular-torsor loop equivalence.}
\upstream{\texttt{Properties.agda}}{ΩEM₁Iso; isConnectedEM₁; isGroupoidEM₁}

\subsection{Why a universe lift appears}
The moduli of small torsors is one universe level higher than an individual small carrier. More precisely, for $G$ at level $\ell$, the checked moduli type is in $\Type_{\ell+1}$. The original $\mathsf{EM}_1(G)$ is at level $\ell$. The local global comparison therefore uses
\begin{equation}
\widetilde{\mathsf{EM}}_1(G)\coloneqq\operatorname{Lift}_{\ell+1}(\mathsf{EM}_1(G)).
\label{eq:lifted-em}
\end{equation}
and the map induced by lowering the lift. The lift is equivalent to the original type, but it places both sides in the same universe required by the Whitehead criterion used below.

\formal{\texttt{UniverseLiftedComparison.agda}}{TorsorDelooping.UniverseLiftedComparison.LiftedEM₁; liftedEm1ToTorsor}{The definition applies the next universe lift to the library delooping, then applies the unlifted comparison map.}

\subsection{A Whitehead criterion}
The upstream theorem $\Omega\mathrm{Equiv}\to\mathrm{Equiv}$ applies to a map between types in the same universe. It asks for an equivalence on set truncations and an equivalence on loop maps at every point. Both sides of \eqref{eq:lifted-em} are connected. Therefore their set truncations are contractible, so the induced map of set truncations is an equivalence. The basepoint loop equivalence propagates to every point by connectedness and elimination into the proposition that a function is an equivalence.

\begin{theorem}[Verified global comparison]
\label{thm:em-torsor}
For every set-level group $G:\mathsf{Group}_\ell$, there is an equivalence of types
\begin{equation}
\widetilde{\mathsf{EM}}_1(G)\simeq \GTorsLevel.
\label{eq:em-equivalence}
\end{equation}
\end{theorem}
\begin{proof}
The lifted source is connected and a $1$-type because lifting preserves these properties and the corresponding statements hold for $\mathsf{EM}_1(G)$. The target is connected and a $1$-type by Section~4. The set-truncation map is an equivalence because both truncations are contractible. The induced loop map is an equivalence at every lifted point by transporting the basepoint result along connectedness. Applying the Whitehead criterion produces the equivalence \eqref{eq:em-equivalence}.
\end{proof}

\formal{\texttt{UniverseLiftedComparison.agda}}{liftedEm1ToTorsorEquiv}{The exact checked result is the lifted source equivalence displayed in Theorem~\ref{thm:em-torsor}.}
\upstream{\texttt{WhiteheadsLemma.agda}}{ΩEquiv→Equiv}

\begin{warning}
Theorem 5.2 is an equivalence of underlying types. The comparison map preserves the specified basepoint by a checked path, but the result is not packaged as a pointed equivalence. Nor does the theorem package a group isomorphism between loop groups. The regular-loop multiplication law used to define the map is the precise group-level compatibility established here.
\end{warning}

\section{The General Path-Torsor Characterization}
\label{ch:path-torsors}

Section~5 compares a particular library delooping with the torsor moduli type. The more intrinsic statement begins with an arbitrary pointed connected $1$-type and asks whether it can be recognized as a torsor moduli type from its loop data.

\subsection{The input data}
\begin{definition}[Pointed loop-group data]
Let $G$ be a set-level group, $\BG: \Type_\ell$, and $pt:\BG$. Pointed loop-group data consist of:
\begin{enumerate}[label=\textup{(\roman*)}]
\item a proof that $\BG$ is connected, written $\operatorname{isConnected}(2,\BG)$;
\item a proof that $\BG$ is a $1$-type, written $\operatorname{isGroupoid}(\BG)$;
\item an equivalence $e:G\simeq(pt=pt)$;
\item a unit path $e(1)=\mathsf{refl}$; and
\item multiplication paths $e(gh)=e(g)\mathbin{\cdot}e(h)$ for all $g,h:G$.
\end{enumerate}
\end{definition}

The last two items are not cosmetic. They orient the chosen loop equivalence so that left multiplication in $G$ acts by left path concatenation. Without them, an underlying equivalence between $G$ and the loop type does not supply the action law used below.

\formal{\texttt{PathTorsor.agda}}{TorsorDelooping.PathTorsor.PointedLoopGroupData}{The five fields are \texttt{isConnectedBG}, \texttt{isGroupoidBG}, \texttt{loopEquiv}, \texttt{loop-unit}, and \texttt{loop-comp}.}

\subsection{Path torsors}
Fix such data and a point $t:BG$. Define the carrier
\begin{equation}
P_t\coloneqq(pt=t).
\end{equation}
Define a left $G$-action by
\begin{equation}
g\cdot q\coloneqq e(g)\mathbin{\cdot}q, \qquad q:pt=t.
\label{eq:path-action}
\end{equation}
Here the first product is group multiplication and the second is path concatenation. The unit law follows from $e(1)=\mathsf{refl}$ and the left unit law for paths. The action law follows from the multiplication path for $e$ and associativity of concatenation:
\begin{align*}
(gh)\cdot q&=e(gh)\mathbin{\cdot}q \\\
&=(e(g)\mathbin{\cdot}e(h))\mathbin{\cdot}q \\\
&=e(g)\mathbin{\cdot}(e(h)\mathbin{\cdot}q)=g\cdot(h\cdot q).
\end{align*}

\formal{\texttt{PathTorsor.agda}}{TorsorDelooping.PathTorsor.pathTorsorAction}{The action is exactly left path concatenation; the code proves its unit and multiplication laws using \texttt{loop-unit} and \texttt{loop-comp}.}

\subsection{The torsor property}
\begin{proposition}
For every $t:BG$, the action in \eqref{eq:path-action} makes $P_t$ into a left $G$-torsor.
\end{proposition}
\begin{proof}
Connectedness gives a merely existing path $p:pt=t$, so $P_t$ is merely inhabited. Fix such a path $p$. The orbit map at $p$ is
\begin{equation}
G\longrightarrow (pt=t), \qquad g\longmapsto e(g)\mathbin{\cdot}p.
\end{equation}
It is the composite of the equivalence $e:G\simeq(pt=pt)$ with right composition by $p$, which is an equivalence from loops at $pt$ to paths from $pt$ to $t$. Therefore every fiber of the orbit map is contractible. By the orbit criterion, the transporter from any $x:P_t$ to any $y:P_t$ is contractible. Together with mere inhabitation, this is the torsor condition.
\end{proof}

\formal{\texttt{PathTorsor.agda}}{TorsorDelooping.PathTorsor.pathTorsorInhabited; pathTorsorOrbitEquiv; pathTorsor}{The inhabitance proof maps connectedness-induced mere paths into propositional truncation; regularity is the fiber contractibility of the composite equivalence.}
\upstream{\texttt{Connected.agda}}{isConnectedSuc→isSurjection; isConnectedPoint}

\subsection{The classifying map and its basepoint}
Sending $t$ to the path torsor $P_t$ defines
\begin{equation}
F:\BG\longrightarrow \GTorsLevel.
\label{eq:path-classifying}
\end{equation}
When $t=pt$, the carrier is the loop type. The chosen equivalence $e:G\simeq(pt=pt)$ is equivariant from the regular torsor to $P_{pt}$ precisely because of the multiplication coherence. The SIP therefore gives a basepoint comparison path
\begin{equation}
b_G=F(pt).
\label{eq:base-comparison}
\end{equation}

\formal{\texttt{PathTorsor.agda}}{TorsorDelooping.PathTorsor.pathTorsorModuli; pathTorsorBaseEquiv; pathTorsorModuli-base}{The local code supplies an equivariant equivalence and then its SIP path.}

\begin{remark}[Basepoint data versus pointed equivalences]
There is also a definitional path stating that the lifted map takes $\operatorname{lift}(pt)$ to $F(pt)$. Combining it with \eqref{eq:base-comparison} produces a basepoint comparison. The final equivalence below is not packaged in Cubical Agda as a pointed equivalence object; these paths are separate checked declarations.
\end{remark}

\formal{\texttt{PathTorsor.agda}}{TorsorDelooping.PathTorsor.liftedPathTorsorModuli-base}{The declaration identifies the lifted basepoint image with the unlifted path-torsor moduli point.}

\subsection{Loop maps}
To compare loops, first observe that a path $q:s=t$ in $\BG$ acts on path torsors by right composition:
\begin{equation}
(pt=s)\longrightarrow(pt=t), \qquad p\longmapsto p\mathbin{\cdot}q.
\end{equation}
This map is equivariant for the left actions, hence SIP turns it into a path $F(s)=F(t)$. Univalence computes the carrier path of that SIP path. This establishes that the action of $F$ on equality paths agrees with the intended transport equivalence.

At the basepoint, equivariant automorphisms of $P_{pt}$ are equivalent to loops at $pt$ by evaluation at $\mathsf{refl}$. Together with the SIP, this proves that the induced basepoint loop map of $F$ is an equivalence. Connectedness then extends this to every point.

\formal{\texttt{PathTorsor.agda}}{TorsorDelooping.PathTorsor.pathTorsorTransportEquiv; pathTorsorModuli-cong; pathTorsorAutLoopIso; isEquiv-pathTorsorModuli-loop}{The calculation uses \texttt{pathFamily-ua}, the SIP carrier computation, and propositional-truncation elimination to move from the basepoint to arbitrary points.}

\subsection{The lifted theorem}
Universe levels again matter: $\BG$ lies in $\Type_\ell$, whereas the torsor moduli type lies one universe higher. Let $\operatorname{Lift}_{\ell+1}(\BG)$ denote the lift. Define the lifted classifying map by lowering the lift and applying $F$.

\begin{theorem}[Universe-lifted path-torsor characterization]
\label{thm:path-torsor}
Given pointed loop-group data for $(G,\BG,pt)$, there is an equivalence of underlying types
\begin{equation}
\operatorname{Lift}_{\ell+1}(\BG)\simeq \GTorsLevel.
\label{eq:path-torsor-equivalence}
\end{equation}
\end{theorem}
\begin{proof}
Lifting preserves connectedness and the $1$-type property. The moduli target is connected and a $1$-type by Section~4. The preceding discussion gives an equivalence on loop maps at all lifted points. Both set truncations are contractible by connectedness, hence the induced map on set truncations is an equivalence. The same Whitehead criterion used in Section~5 yields \eqref{eq:path-torsor-equivalence}.
\end{proof}

\formal{\texttt{PathTorsor.agda}}{TorsorDelooping.PathTorsor.liftedPathTorsorModuliEquiv}{The exact checked result is the universe-lifted equivalence displayed in Theorem~\ref{thm:path-torsor}.}

\subsection{Relation to Lemma 5.8}
Equation \eqref{eq:path-torsor-equivalence} is a universe-lifted path-torsor characterization associated with Lemma 5.8 of Myers--Sati--Schreiber. It is not an unconditional theorem asserting $\BG\simeq\GTors$ from a bare informal phrase ``the loop group is $G$.'' The input explicitly includes connectedness, $1$-truncation, a chosen loop-type equivalence, and chosen unit and multiplication coherences. It also does not identify the result with a geometric principal-bundle construction.

\section{Mathematical Perspective}
\label{ch:perspective}

\subsection{What has been established}
The development establishes a generic theorem about set-level groups and torsors in Cubical Agda.
\begin{enumerate}[label=\arabic*.]
\item Left $G$-actions, left $G$-torsors, equivariant maps, and the regular torsor are defined and verified.
\item The moduli type of left torsors is constructed by a univalent structured-type construction. Its paths are equivalent to equivariant equivalences.
\item The torsor moduli is connected and a $1$-type. Its loop type at the regular torsor is equivalent, as a type, to the carrier of $G$.
\item The restricted coherence for regular right translations is proved: the loops used in the $\mathsf{EM}_1$ recursor respect multiplication in the required orientation.
\item The lifted library delooping is equivalent, as a type, to the torsor moduli.
\item More generally, a pointed connected $1$-type equipped with explicit compatible loop-group data admits the universe-lifted path-torsor equivalence of Section~6.
\end{enumerate}

These results give a precise internal realization of the familiar slogan that a connected delooping is the moduli of torsors for its loop group. The qualifier ``precise'' is important: the universe lift and every compatibility datum are present in the formal theorem.

\subsection{Local results and inherited results}
The local work proves the torsor-specific constructions: the transporter formulation, regularity, the moduli structure, connectedness, the regular-loop coherence, and the two global comparison maps. It relies on existing Cubical Agda mathematics for univalence, SIP, higher inductive Eilenberg--Mac Lane types, truncations, connectedness tools, and the Whitehead criterion.

One should distinguish this division of labour from a distinction in mathematical validity. The local proof terms are checked relative to the upstream library. The library results are also checked in the pinned checkout. Appendix B records the key interfaces and the exact pinned revision.

\subsection{Three levels of structure}
It is useful to separate three statements that are easily conflated.
\begin{enumerate}[label=\textbf{Level \arabic*:},leftmargin=2.3cm]
\item An equivalence of underlying types, such as \eqref{eq:em-equivalence} or \eqref{eq:path-torsor-equivalence}. This is what the two final declarations provide.
\item A basepoint comparison path. For the $\mathsf{EM}_1$ map, the image of the basepoint is definitionally the regular torsor. For the general path-torsor map, the source basepoint is related to the regular torsor through the separately checked SIP path of Equation \eqref{eq:base-comparison}.
\item A group-level identification. Such a statement would require a group structure on the loop type and a packaged proof that a specified equivalence preserves multiplication, unit, and inverse. The current milestone proves the regular multiplication coherence required by the recursor, but does not expose a new group-isomorphism record as its endpoint.
\end{enumerate}

\subsection{Limitations}
The text has deliberately stayed at the generic level. It does not formalize the complex multiplicative group $\mathbb{C}^{\times}$, a topology on that group, smooth or geometric principal torsors, or a comparison with a conventional construction of $B\mathbb{C}^{\times}$. It constructs no dependent type family over a braid-group or torsor delooping, no monodromy representation, and no computational transport operation.

Accordingly, this milestone is not a formalization of Definition 6.7 of Myers--Sati--Schreiber. The eventual programme needs additional ingredients: a suitable model of the relevant group or higher group, a family of types over its classifying type, a proof that the desired loops act through transport, and a computational analysis of that transport.

\subsection{A route toward later applications}
The generic result is nevertheless useful for later work. If a candidate group $G$ and a candidate pointed connected $1$-type $\BG$ are supplied, Section~6 says exactly which data are needed to compare $\BG$ with the moduli of $G$-torsors. In particular, an informal identification of loops is insufficient: one must expose its unit and multiplication coherences. This is a valuable design constraint for any later construction involving $\mathbb{C}^{\times}$, torsor classifiers, or braid-group actions.

\begin{remark}
The same discipline applies to computational claims. A type-theoretic existence proof of a transport map is not yet evidence that the transport has useful reduction behaviour. That question remains open and should be studied with concrete dependent families rather than inferred from the present delooping theorem.
\end{remark}

\subsection{Reading the formalization notes}
Every correspondence note in this article names a local source-relative path and a checked declaration. Such a note means that the indicated construction, or the indicated exact statement, occurs in the local Cubical Agda source. It does not mean that the prose proof is definitionally identical to the proof term, nor that every intermediate sentence has a separate declaration. Upstream notes identify imported verified interfaces. Appendix A is the consolidated statement-to-declaration index.

\appendix
\section{Statement-to-Declaration Index}
\label{app:index}

This is a provenance index, not a second proof. ``Local'' means a declaration in \texttt{src/TorsorDelooping/}; ``upstream'' means the pinned Cubical checkout. Repository-relative paths can later be attached to an immutable milestone tag without changing the mathematical text.

\subsection*{Local construction index}
\begingroup\small\raggedright
\begin{description}[font=\normalfont\bfseries]
\item[Actions and torsors]
\texttt{src/TorsorDelooping/LeftAction.agda}\par
\agda{TorsorDelooping.LeftAction.LeftAction}, \agda{TorsorDelooping.LeftAction.Equivariant}.\par
\texttt{src/TorsorDelooping/LeftTorsor.agda}\par
\agda{TorsorDelooping.LeftTorsor.IsLeftTorsor}.\par
\texttt{src/TorsorDelooping/Regular.agda}\par
\agda{regularLeftAction}, \agda{regularLeftTorsor}.

\item[Equivariant maps]
\texttt{src/TorsorDelooping/Morphisms.agda}\par
\agda{equivariant-id}, \agda{equivariant-comp}, \agda{orbitEquiv}, \agda{equivariantMapIsEquiv}.

\item[Moduli and structure identity]
\texttt{src/TorsorDelooping/Moduli.agda}\par
\agda{LeftTorsorModuli}, \agda{TorsorEquiv}, \agda{torsorPathEquiv}.

\item[Structured-path embedding]
\texttt{src/TorsorDelooping/SIPPathEmbedding.agda}\par
\agda{caracTorsor≡}. This proves that structured torsor paths are determined by their carrier paths; it is used in the SIP coherence calculation.

\item[Delooping structure]
\texttt{src/TorsorDelooping/Delooping.agda}\par
\agda{torsorBase}, \agda{isConnectedLeftTorsorModuli}, \agda{isGroupoidLeftTorsorModuli}, \agda{regularAutIsoGroup}, \agda{rightRegularAut}, \agda{torsorLoopEquivGroup}.

\item[Composition coherence]
\texttt{src/TorsorDelooping/LoopGroup.agda}: \agda{rightRegularAut-comp}.\par
\texttt{src/TorsorDelooping/SIPCoherence.agda}: \agda{torsorPath-carrier}, \agda{regularTorsorLoop}, \agda{regularTorsorLoop-comp}.

\item[Comparison with the library delooping]
\texttt{src/TorsorDelooping/EM1Comparison.agda}: \agda{em1ToTorsor}, \agda{isEquiv-em1ToTorsor-baseLoop}.\par
\texttt{src/TorsorDelooping/UniverseLiftedComparison.agda}: \agda{LiftedEM₁}, \agda{liftedEm1ToTorsor}, \agda{liftedEm1ToTorsorEquiv}.

\item[General path-torsor theorem]
\texttt{src/TorsorDelooping/PathTorsor.agda}\par
\agda{PointedLoopGroupData}, \agda{pathTorsorAction}, \agda{pathTorsor}, \agda{pathTorsorModuli}, \agda{pathTorsorModuli-base}, \agda{liftedPathTorsorModuli-base}, \agda{liftedPathTorsorModuliEquiv}.
\end{description}
\endgroup

\subsection*{Narrative statements without a one-to-one declaration}
The elementary examples of affine spaces and phase torsors, and the expository proofs supplied in Sections~1--4, are mathematical exposition. They explain the checked definitions and constructions but are not separate Agda modules. Likewise, the discussion of later applications to $\mathbb{C}^{\times}$ and Definition~6.7 is explicitly prospective, not a verified construction.

\section{Upstream Provenance and Reproducibility}
\label{app:provenance}

\subsection{Pinned environment}
The audited environment is Agda \texttt{2.8.0-3d04bac} with the project-local Cubical Agda checkout at commit
\begin{center}
\texttt{87af036ebd1c87956cf08fea44e62c7cf76392ca}.
\end{center}
The project library descriptor enables Cubical and guardedness options. The local development is checked module by module with explicit include paths, rather than relying on a globally registered Agda library.

\subsection{Reproducibility command}
From the repository root, the complete local check is:
{\small\begin{verbatim}
for file in $(find src/TorsorDelooping -name '*.agda' | sort); do
  agda --safe -i src -i cubical "$file" || exit $?
done
\end{verbatim}}
At the time of this draft the command ranges over twelve modules. A successful run returns exit status zero and no Agda diagnostics. The checked modules use neither postulates nor unresolved holes. The option \texttt{--lossy-unification} appears in the universe-lifted and path-torsor modules as a unification heuristic; the modules nevertheless pass Agda's \texttt{--safe} check and it does not add an axiom or an unsafe option.

\subsection{Principal upstream interfaces}
\begingroup\small\raggedright
\begin{description}[font=\normalfont\bfseries]
\item[Group right translations]
\texttt{cubical/Cubical/Algebra/Group/Properties.agda}\par
\agda{·GroupAutomorphismR}.
\item[Univalence composition]
\texttt{cubical/Cubical/Foundations/Univalence.agda}\par
\agda{uaCompEquiv}.
\item[Structure Identity Principle]
\texttt{cubical/Cubical/Foundations/SIP.agda}\par
\agda{SIP}.
\item[The $\mathsf{EM}_1$ recursor]
\texttt{cubical/Cubical/HITs/EilenbergMacLane1/Properties.agda}\par
\agda{rec'}.
\item[Connectedness and loop calculation for $\mathsf{EM}_1$]
\texttt{cubical/Cubical/Homotopy/EilenbergMacLane/Properties.agda}\par
\agda{isConnectedEM₁}, \agda{isGroupoidEM₁}, \agda{ΩEM₁Iso}.
\item[Whitehead criterion]
\texttt{cubical/Cubical/Homotopy/WhiteheadsLemma.agda}\par
\agda{ΩEquiv→Equiv}.
\item[Connectedness tools]
\texttt{cubical/Cubical/Homotopy/Connected.agda}\par
\agda{isConnectedSuc→isSurjection}, \agda{isConnectedRetractFromIso}.
\end{description}
\endgroup

The SIP and univalence interfaces are imported rather than reconstructed locally. The local contribution is the construction of the relevant univalent torsor structure and the concrete coherence calculation allowing the $\mathsf{EM}_1$ recursor to be used.

\subsection{Repository navigation}
The source tree is organized in dependency order:
\begin{verbatim}
LeftAction -> LeftTorsor -> Regular, Morphisms -> Moduli
Moduli -> Delooping, LoopGroup -> SIPCoherence
                                      -> EM1Comparison
EM1Comparison -> UniverseLiftedComparison -> PathTorsor
\end{verbatim}
The repository's \texttt{alignment/records.jsonl} contains machine-readable statement-to-source records. The source-relative paths in Appendix~A are deliberately retained instead of URLs to a moving branch. Once this milestone is committed and tagged, they can be linked to the corresponding immutable tag without altering the claims of this article.

\subsection{Boundary of verification}
The local code proves the type-theoretic statements indexed in Appendix~A. The paper's elementary mathematical explanations are not individually formal artifacts, and its bibliographic comparison with Lemma~5.8 is interpretive. No claim in this paper establishes a geometric model for $B\mathbb{C}^{\times}$, a computational transport operation, a braid-group action, or Definition~6.7. These remain future work.

\begin{thebibliography}{9}
\bibitem{cubical} Cubical Agda contributors. \emph{Cubical Agda standard library}. Pinned project checkout at the revision above. \url{https://github.com/agda/cubical}.
\bibitem{agda} The Agda Team. \emph{Agda 2.8.0-3d04bac}. \url{https://agda.readthedocs.io/}.
\bibitem{mss} David Jaz Myers, Hisham Sati, and Urs Schreiber. \emph{Topological Quantum Gates in Homotopy Type Theory}. \emph{Communications in Mathematical Physics} \textbf{405}, Article 172 (2024). \url{https://doi.org/10.1007/s00220-024-05020-8}. arXiv:\url{https://arxiv.org/abs/2303.02382}.
\end{thebibliography}

\end{document}
